\section{Removing the vertex costs}
\label{sec:unweighting}

The weighted gap converts to an unweighted gap without requiring a common
denominator for all vertex costs.  Rounding upward at a scale determined
by the number of vertices controls the additive loss. Independent-copy
blowups also occur in earlier DFVS hardness reductions
\cite[Section~2]{GuruswamiLee2016}. The rounding below allows arbitrary
rational costs of polynomial bit length without expanding a common
denominator.

\begin{lemma}[Rounding and cloning]
\label{lem:round-and-clone}
Let $G$ be a loopless digraph on $S\ge1$ vertices with positive rational
vertex costs $w_v\le C$.  Set $c_v=\lceil Sw_v\rceil$.  Replace every vertex
$v$ by an independent set $\mathcal C_v$ of $c_v$ unweighted vertices, and
replace every arc $u\to v$ by all arcs from $\mathcal C_u$ to
$\mathcal C_v$.  Call the resulting digraph $\widehat G$.  Then
\begin{equation}
  \DFVS(\widehat G)
  =\min\left\{\sum_{v\in F}c_v:
      F\text{ is a feedback vertex set of }G\right\}.
  \label{eq:clone-exact-optimum}
\end{equation}
The number of vertices of $\widehat G$ is at most $CS^2+S$.  For fixed
$C$, this is a deterministic polynomial construction in the encoding size
of $G$ and its costs.  Moreover:
\begin{enumerate}[label=(\roman*)]
  \item a feedback vertex set of $G$ of cost at most $a$ yields a feedback
  vertex set of $\widehat G$ of size at most $S(a+1)$;
  \item if every feedback vertex set of $G$ has cost greater than $b$,
  then $\DFVS(\widehat G)>Sb$.
\end{enumerate}
\end{lemma}

\begin{proof}
Let $F$ be a feedback vertex set of $G$.  Deleting every cluster
$\mathcal C_v$ with $v\in F$ leaves only clusters over the acyclic digraph
$G-F$.  Lift a topological order of $G-F$ by placing all vertices of each
cluster consecutively.  Since clusters are independent, this is a
topological order of the surviving vertices of $\widehat G$.  It gives a
feedback vertex set of size $\sum_{v\in F}c_v$.

Conversely, let $Q$ be any feedback vertex set of $\widehat G$, and define
\[
  F_Q=\{v\in V(G):\mathcal C_v\subseteq Q\}.
\]
If $G-F_Q$ contained a directed cycle, every cluster over that cycle would
contain a surviving vertex.  Choosing one such vertex from each cluster
would lift the original cycle to a directed cycle of $\widehat G-Q$,
because all the intercluster arcs in the requisite directions are present.
This is impossible.  Thus $F_Q$ is a feedback vertex set of $G$, and
\[
  \sum_{v\in F_Q}c_v\le\abs{Q}.
\]
Together with the preceding construction, this proves
Equation~\eqref{eq:clone-exact-optimum}.  In particular, partial deletions
from a cluster cannot improve the optimum beyond the corresponding
choice of completely deleted clusters.

For every $v$ we have
\[
  Sw_v\le c_v<Sw_v+1.
\]
Summing over any feedback vertex set $F$ gives
\[
  S\sum_{v\in F}w_v
  \le\sum_{v\in F}c_v
  \le S\sum_{v\in F}w_v+S.
\]
This proves both gap implications by
Equation~\eqref{eq:clone-exact-optimum}.  Equivalently, the rounded costs
$c_v/S$ increase the cost of any deletion set by at most one.

Finally,
\[
  \abs{V(\widehat G)}=\sum_vc_v\le CS^2+S.
\]
There are at most the square of this number of arcs.  Each $c_v$ is
computed exactly by integer arithmetic on the rational input cost.
For fixed $C$, the vertices and arcs can therefore be listed in
polynomial time.  Positivity of the costs ensures that every cluster is
nonempty before deletions.  The assumption that $G$ is loopless ensures
that independent clusters realize precisely the prescribed construction.
\end{proof}

\begin{remark}[Avoiding opposite arcs]
\label{rem:no-opposite-arcs}
The construction above already gives loopless digraphs, and unrestricted
DFVS permits opposite arcs.  If opposite arcs are to be excluded, subdivide
every weighted arc using a distinct private vertex of cost $D+1$ before
applying Lemma~\ref{lem:round-and-clone}.  A feedback vertex set of cost at
most $4$ in the original graph still works in the subdivided graph: a
cycle avoiding that set would project to a directed closed walk, hence
to a directed cycle, in the original remaining graph.  On the other hand,
a deletion set of cost at most $D$ cannot contain any subdivision vertex.
Every surviving original cycle would then lift through its private
subdivisions, so the weighted soundness bound is preserved as well.
Subdivision of every arc of a loopless digraph leaves no pair of opposite
arcs, and the subsequent independent-cluster construction preserves this
property.  The additional vertices have constant cost and polynomial
number, so the size bounds needed for rounding remain valid.
\end{remark}

\begin{proof}[Proof of Theorem~\ref{thm:main}]
It suffices to prove the claim for every fixed integer $A\ge1$: given any
fixed real factor, choose an integer at least as large.  For this integer
$A$, choose the parameters of Section~\ref{sec:parameters} and start from
the NP-hard promise problem of Theorem~\ref{thm:games} at the resulting
fixed $\theta$.

By Proposition~\ref{prop:polynomial}, the weighted reduction is
deterministic and polynomial in the game input size, and all its costs are
bounded by the constant $C_A$.  Proposition~\ref{prop:completeness} gives
a feedback vertex set of cost at most $4$ in the YES case, while
Proposition~\ref{prop:weighted-soundness} shows that every feedback vertex
set has cost greater than $D=10A$ in the NO case.  Let $S$ be the number of
weighted vertices.  Lemma~\ref{lem:round-and-clone} gives an unweighted
digraph $\widehat G$ with the gap
\[
  \begin{array}{ll}
    \text{YES:}&\DFVS(\widehat G)\le5S,\\[2pt]
    \text{NO:}&\DFVS(\widehat G)>DS.
  \end{array}
\]
This last transformation is polynomial because $C_A$ is fixed.  Set
$k=5S$ and output the pair $(\widehat G,k)$.  In the YES case,
$\DFVS(\widehat G)\le k$, whereas in the NO case
\[
  \DFVS(\widehat G)>DS=10AS>5AS=Ak.
\]
Thus this pair has the promised gap of Theorem~\ref{thm:main}.

A deterministic polynomial-time $A$-approximation would return a feasible
feedback vertex set of size at most $5AS<DS$ on every YES instance.  On a
NO instance, every feasible feedback vertex set has size greater than
$DS$.  Comparing the returned size with the known threshold $DS$ would
therefore distinguish the NP-hard promise cases in polynomial time.
This proves the asserted NP-hardness for every fixed constant factor,
and its deterministic consequence under $\mathrm P\ne\mathrm{NP}$.
\end{proof}


A \emph{directed feedback arc set} is a set of arcs whose deletion leaves
an acyclic digraph.  Write $\operatorname{DFAS}(H)$ for its minimum
cardinality.

\begin{corollary}[Directed feedback arc set]
\label{cor:feedback-arc}
For every fixed real $A\ge1$, approximating directed feedback arc set
within factor $A$ is NP-hard on unweighted simple loopless digraphs.
More precisely, there is a deterministic polynomial reduction producing
such a digraph $H$ and a positive integer $k$ with the promise gap
\[
  \begin{array}{ll}
    \text{YES:}&\operatorname{DFAS}(H)\le k,\\[2pt]
    \text{NO:}&\operatorname{DFAS}(H)>Ak.
  \end{array}
\]
\end{corollary}

\begin{proof}
Let $G=(V,E)$ be an unweighted loopless digraph, with $n=\abs{V}$ and
$m=\abs{E}$.  Split each vertex $v$ into $v^-,v^+$ and insert the arc
$a_v:v^-\to v^+$.  For every original arc $e=(u,v)$, introduce $n+1$
fresh private vertices $w_{e,j}$ and the paths
\[
  u^+\longrightarrow w_{e,j}\longrightarrow v^-\qquad(1\le j\le n+1).
\]
Call the resulting unweighted digraph $H$.  All private vertices are
distinct, so $H$ has no loops or repeated arcs, even if $G$ has opposite
arcs.  It has $2n+(n+1)m$ vertices and $n+2(n+1)m$ arcs and is explicitly
constructible in polynomial time.

If $S$ is a feedback vertex set of $G$, deleting the arcs $a_v$ for
$v\in S$ destroys every directed cycle of $H$: a surviving cycle would
alternate distinguished arcs with connector paths and project to a
directed closed walk in $G-S$.  Hence
$\operatorname{DFAS}(H)\le\DFVS(G)\le n$.

Conversely, let $F$ be any feedback arc set of $H$.  If $\abs{F}\ge n$,
return $S=V$, which is a feedback vertex set of size at most $\abs{F}$.
If $\abs{F}<n$, return $S=\{v:a_v\in F\}$, again of size at most
$\abs{F}$.  Were there a directed cycle in $G-S$, every distinguished arc
needed to lift it would survive.  For each original arc of the cycle,
its $n+1$ connector paths are pairwise edge-disjoint, so at least one
survives deletion of $F$.  Choosing these paths would lift the cycle to
$H-F$, a contradiction.  Thus the returned $S$ is always a feedback
vertex set, and this recovery is polynomial time.

The two maps give $\operatorname{DFAS}(H)=\DFVS(G)$.  They also preserve
approximation solutions: an $A$-approximate $F$ yields
$\abs{S}\le\abs{F}\le A\operatorname{DFAS}(H)=A\DFVS(G)$.
Composing this construction with Theorem~\ref{thm:main} and retaining its
positive integer threshold $k$ proves the stated gap.
\end{proof}
